% ============================================================================
% templates/report/silabo/silabo.tex — plantilla: sílabo
% ----------------------------------------------------------------------------
% Sílabo de curso. SIN diseño: todo lo pone academic-report + la identidad.
% Uso: scripts/build.sh silabo.tex   (LuaLaTeX)
% ============================================================================
\documentclass{academic-report}
\tipodocumento{Sílabo}
\universidad{Universidad Nacional de San Cristóbal de Huamanga}
\facultad{Facultad de Ciencias Económicas, Administrativas y Contables}
\escuela{Escuela Profesional de Economía}
\curso{Nombre del curso}\codigocurso{XX-000}
\docente{Edison Achalma B.Sc. Econ.}\periodo{Semestre 2026-I}

\begin{document}
\encabezado

\section{Datos generales}
\datofila{Créditos}{4}
\datofila{Horas semanales}{Teoría 3 · Práctica 2}
\datofila{Prerrequisito}{Curso previo}
\datofila{Ciclo}{V}

\section{Sumilla}
Descripción general del curso: naturaleza, propósito y aporte al perfil del egresado.

\section{Competencias}
\begin{itemize}
  \item Competencia 1
  \item Competencia 2
\end{itemize}

\section{Resultados de aprendizaje}
\begin{enumerate}
  \item Resultado 1
  \item Resultado 2
\end{enumerate}

\section{Unidades didácticas}
\begin{description}
  \item[Unidad 1.] Título — contenidos principales.
  \item[Unidad 2.] Título — contenidos principales.
\end{description}

\section{Metodología}
Estrategias de enseñanza-aprendizaje, recursos y modalidad.

\section{Evaluación}
\begin{table}[h]\centering
\begin{tabularx}{\linewidth}{Xc}
\toprule \textbf{Instrumento} & \textbf{Peso} \\\midrule
Prácticas calificadas & 30\,\% \\
Examen parcial & 30\,\% \\
Examen final & 40\,\% \\\bottomrule
\end{tabularx}
\end{table}

\section{Cronograma}
Ver el calendario del curso.

\section{Bibliografía}
\begin{itemize}
  \item Autor (Año). \emph{Título}. Editorial.
\end{itemize}

\end{document}
